\section{Discussion}

\label{sec:discussion}

\subsection{Advantages of Temporal Sequence Processing}

The TSN architecture's primary advantage over bi-temporal methods is robustness to confounding factors.
Seasonal phenological cycles, which cause the majority of false positives in NDVI-based approaches, are modeled explicitly by the temporal attention module.
Cloud contamination, which delays detection in single-date systems, is handled by attending to cloud-free timesteps within the sequence.
Gradual degradation, invisible in any single image pair, becomes detectable through cumulative temporal signal analysis.

\subsection{Limitations}

Several limitations constrain the current system.
First, the 10-meter spatial resolution of Sentinel-2 limits detection of narrow linear features (roads, trails) and small-scale selective logging.
Integration with high-resolution commercial imagery (Planet, Maxar) for targeted areas would address this.
Second, the TSN model is trained on labeled change events, which are biased toward conspicuous changes; subtle degradation may be underrepresented.
Third, the HFI is species-agnostic; extending it to species-specific connectivity metrics requires integration with species distribution models.

\subsection{Scalability}

The current deployment covers 2.4 million km$^2$.
Scaling to global coverage (approximately 50 million km$^2$ of non-ice terrestrial surface) would require approximately 20x the current computational resources, bringing the cost to approximately \$135,000/month.
This is within the budget range of major conservation organizations and could be supported through the Zoo Network's decentralized compute infrastructure.

\subsection{Integration with Zoo Ecosystem}

ZooSat integrates with the broader Zoo ecosystem:
\begin{itemize}[leftmargin=*]
    \item Change detection outputs are stored in the Zoo Data Commons (ZIP-009) with full provenance tracking.
    \item HFI trajectories are published through the ZooCoord research coordination protocol (ZIP-008).
    \item Alert delivery uses the Zoo Network's decentralized messaging infrastructure.
    \item Model training leverages ZooFed's federated learning framework for cross-site improvement.
\end{itemize}

% ============================================================
% 9. Conclusion
% ============================================================
